\section{Arbitrary-index Stieltjes pairings}
\label{sec:stieltjes}

\subsection{The complete coefficient theorem}
Expansion of the subtracted term in \eqref{eq:Z-definition-intro} gives
\begin{equation}\label{eq:Z-G-expansion}
 Z_{1+u}(z)=\sum_{m\ge0}\frac{(-1)^m}{m!}G_m(z)u^m.
\end{equation}
For $\Rea a,\Rea b>0$, define
\begin{equation}\label{eq:Imn}
 I_{m,n}(A)=\frac1{2\pi}\int_{\R}
            G_m(a+\ii y)G_n(b-\ii y)\dd y,
 \qquad A=a+b.
\end{equation}
The following theorem both proves that this depends only on $A$ and
specifies a finite answer at every pair of indices.

\begin{theorem}[All-index Stieltjes closure]\label{thm:all-stieltjes}
For all integers $m,n\ge0$, the integral \eqref{eq:Imn} converges absolutely
and
\begin{equation}\label{eq:all-jets}
 \boxed{\displaystyle
 I_{m,n}(A)=(-1)^{m+n}m!n!\,[u^mv^n]
      C(u,v)K(1+u+v,A),}
\end{equation}
where
\begin{align}
 C(u,v)&=\frac{\Gamma(1+u+v)}{\Gamma(1+u)\Gamma(1+v)},\label{eq:C}\\
 \log C(u,v)&=\sum_{r\ge2}\frac{(-1)^r\zeta(r)}r
          \{(u+v)^r-u^r-v^r\}.\label{eq:log-C}
\end{align}
Write $K(1+\epsilon,A)=\sum_{\ell\ge0}k_\ell(A)\epsilon^\ell$.
Its coefficients are
\begin{align}
 k_\ell(A)={}&\frac{\zeta^{(\ell)}(0,A)}{\ell!}
  -\frac{2\zeta^{(\ell+1)}(0,A)}{(\ell+1)!}
  +(1-A)\frac{(-1)^\ell\gamma_\ell(A)}{\ell!}\notag\\
 &-A\sum_{j=0}^{\ell+1}\frac{(-\log A)^j}{j!}.
       \label{eq:k-coefficients}
\end{align}
In particular, $I_{m,n}$ belongs to the finite algebra generated by
$A$, $\log A$, the constants $\zeta(2),\ldots,\zeta(m+n)$,
$\gamma_\ell(A)$ for $\ell\le m+n$, and
$\zeta^{(j)}(0,A)$ for $j\le m+n+1$.
\end{theorem}
\begin{proof}
The map $u\mapsto Z_{1+u}(a+\ii y)$ is holomorphic as an $L^2(\R)$ vector
in a neighborhood of zero, by the estimates preceding \cref{thm:basic}.
Its derivatives are precisely $(-1)^mG_m$. Hence every $G_m$ belongs to
$L^2$, and the product in \eqref{eq:Imn} is absolutely integrable.
Taking the two Taylor coefficients in \eqref{eq:basic-pair} gives
\eqref{eq:all-jets}. The Taylor expansion of $\log\Gamma(1+u)$ proves
\eqref{eq:log-C}; its linear Euler-constant terms cancel.

For the explicit coefficients, write
\[
 K(1+\epsilon,A)=\left(1-\frac2\epsilon\right)\zeta(\epsilon,A)
 +(1-A)\zeta(1+\epsilon,A)
 +\frac{A e^{-\epsilon\log A}}{\epsilon(\epsilon-1)}.
\]
The pole coefficient vanishes as already checked. Since
\[
 \frac{Ae^{-\epsilon\log A}}{\epsilon(\epsilon-1)}
 =-A\sum_{r\ge0}\epsilon^{r-1}
          \sum_{j=0}^r\frac{(-\log A)^j}{j!},
\]
its coefficient of $\epsilon^\ell$ is the last term of
\eqref{eq:k-coefficients}. The other terms follow from the Taylor and
Laurent series of the two zeta factors. Only total degree at most $m+n$
of $C$ is used, proving the finite-alphabet assertion.
\end{proof}

This theorem is an actual finite recipe. If
$C(u,v)=\sum c_{r,k}u^rv^k$, then equivalently
\begin{equation}\label{eq:finite-jet-recipe}
 I_{m,n}(A)=(-1)^{m+n}m!n!
  \sum_{r=0}^m\sum_{k=0}^n c_{r,k}
       \binom{m+n-r-k}{m-r}k_{m+n-r-k}(A).
\end{equation}
No two-variable special function or undetermined integration constant
remains. This is a closure statement, not a claim that all its generators
are algebraically independent or reducible to elementary constants.

\subsection{Low-order identities and a useful degeneration}
Let $\mu_j(A)=\partial_w^jK(1,A)=j!k_j(A)$. Direct coefficient extraction
from \eqref{eq:log-C} gives
\begin{align}
 I_{0,0}&=\mu_0,&
 I_{m,0}&=(-1)^m\mu_m,\label{eq:jet-row0}\\
 I_{1,1}&=\mu_2+\zeta(2)\mu_0,\label{eq:jet11}\\
 I_{2,1}&=-\mu_3-2\zeta(2)\mu_1+2\zeta(3)\mu_0,\label{eq:jet21}\\
 I_{3,1}&=\mu_4+3\zeta(2)\mu_2-6\zeta(3)\mu_1
                   +6\zeta(4)\mu_0,\label{eq:jet31}\\
 I_{2,2}&=\mu_4+4\zeta(2)\mu_2-8\zeta(3)\mu_1
                   +\{6\zeta(4)+2\zeta(2)^2\}\mu_0.
                   \label{eq:jet22}
\end{align}
Here and below, every $\mu_j$ in one equation has the same argument $A$.
For example, the coefficient of $u^2v^2$ in $C$ is
$\frac32\zeta(4)+\frac12\zeta(2)^2$; the factor $2!2!$ explains the last
coefficient in \eqref{eq:jet22}. This factorial is easy to lose if ordinary
and exponential generating functions are mixed.

At $A=1$, the entire explicit Stieltjes term in
\eqref{eq:k-coefficients} disappears:
\begin{equation}\label{eq:A-one-degeneration}
 k_\ell(1)=\frac{\zeta^{(\ell)}(0)}{\ell!}
          -\frac{2\zeta^{(\ell+1)}(0)}{(\ell+1)!}-1.
\end{equation}
Thus every vertical pairing with $a+b=1$ is expressible using Riemann zeta
order derivatives at zero and ordinary positive zeta values. This does
not say that the individual Stieltjes functions disappear pointwise;
the simplification occurs after pairing.

\subsection{A second, kernel-polynomial description}
Define real-coefficient polynomials $E_m(L)$ by
\begin{equation}\label{eq:E-generating}
 \sum_{m\ge0}E_m(L)\frac{v^m}{m!}
          =\frac{e^{-vL}}{\Gamma(1-v)}.
\end{equation}
For example,
\[
 E_0=1,\qquad E_1=-(L+\gamma),\qquad
 E_2=(L+\gamma)^2-\zeta(2).
\]
Differentiating \eqref{eq:Z-laplace} gives
\begin{align}
 G_m(z)&=\int_0^\infty e^{-zx}h(x)E_m(\log x)\dd x,
               \label{eq:G-bell-laplace}\\
 I_{m,n}(A)&=\int_0^\infty e^{-Ax}h(x)^2
                  E_m(\log x)E_n(\log x)\dd x.
               \label{eq:Stieltjes-bell-pair}
\end{align}
All these differentiations are justified by the same compact-parameter
bounds used in the proof of \cref{thm:all-stieltjes}.
Equations~\eqref{eq:all-jets} and \eqref{eq:Stieltjes-bell-pair} provide two
useful computational routes: the former uses a finite zeta/Stieltjes
alphabet, while the latter uses one rapidly decaying quadrature with a kernel regular for positive $x$. They are independently implemented in the verification code.

The equal-scale identity also supplies all derivatives in the horizontal
arguments: one differentiates \eqref{eq:all-jets} in $A$, or inserts powers
of $-x$ in \eqref{eq:Stieltjes-bell-pair}. No new spectral-order restriction
arises from a nonnegative number of such derivatives.

\subsection{A complex-safe Hermite identity}
The generalized Stieltjes parameter must be treated holomorphically, not by
applying a real-part formula outside its real domain. A useful exact
representation, valid at every index, is the following.

\begin{proposition}[Holomorphic Hermite representation]\label{thm:complex-Hermite}
For $\Rea z>0$ and every integer $m\ge0$,
\begin{align}
 G_m(z)={}&\frac{\log^m z}{2z}
  +\ii\int_0^\infty\frac1{e^{2\pi t}-1}
       \left\{\frac{\log^m(z+\ii t)}{z+\ii t}
                  -\frac{\log^m(z-\ii t)}{z-\ii t}\right\}\dd t.
                     \label{eq:complex-Hermite}
\end{align}
The integral is absolutely convergent and holomorphic in $z$ throughout
the right half-plane. For nonreal $z$, neither term may be replaced by
the complex conjugate of the other.
\end{proposition}
\begin{proof}
The standard Hermite representation of Hurwitz zeta
\cite{dlmf-zeta}, after subtracting its spatial principal term, is
\[
 Z_s(z)=\frac1{2z^s}
   +\ii\int_0^\infty
        \frac{(z+\ii t)^{-s}-(z-\ii t)^{-s}}{e^{2\pi t}-1}\dd t.
\]
It can equivalently be obtained from \eqref{eq:Z-laplace} using
$h(x)=1/2+2\int_0^\infty\sin(tx)/(e^{2\pi t}-1)\dd t$.
The two powers have a difference $O(t)$ at zero, uniformly on compact
subsets of $s\in\C$, $\Rea z>0$; the denominator also has a simple zero.
At infinity the exponential denominator dominates all powers and
logarithms arising from any fixed spectral derivative. Therefore the
combined integral is jointly holomorphic and may be differentiated in
$s$ at $s=1$. Multiplying the $m$th derivative by $(-1)^m$ proves
\eqref{eq:complex-Hermite}. Holomorphic continuation in $z$ from positive
real values gives the entire right half-plane; the paired expression,
not a real-part projection, is what makes this continuation legitimate.
\end{proof}

This representation is classical in origin, but spelling out its paired
complex form is important for the verification of the present vertical-line
identities. \Cref{sec:software-audit} records a concrete numerical failure
when a real-part implementation is used at nonreal $z$, and supplies both
\eqref{eq:complex-Hermite} and Cauchy-coefficient extraction as repairs.
